\begin{register}{H}{tselect}{0x7A0}\label{regdesc:LPTSELECT}
\regfield{tselect}{32}{0}{00000000000000000000000000000000}%
\reglabel{Reset}\regnewline%
\begin{regdesc}\begin{reglist}
\label{fielddesc:LPTSELECT}\item [tselect] 配置选择触发器单元。(R/W)
\end{reglist}\end{regdesc}
\end{register}

%%%%%%%%%
\begin{register}{H}{tdata1}{0x7A1}\label{regdesc:LPTDATA1}
\regfield{type}{4}{28}{0010}%
\regfield{dmode}{1}{27}{1}%
\regfield{data}{27}{0}{0x1040}%
\reglabel{Reset}\regnewline%
\begin{regdesc}\begin{reglist}
\label{fielddesc:LPTDATA1TYPE}\item [type] 表示触发器类型。仅支持匹配类型 (0x2)，此域保留。(RO)
\label{fielddesc:LPTDATA1DMODE}\item [dmode] 如果某触发器正在被调试器使用，则此域置为 1。仅支持调试模式使用，此域保留。(RO)
\label{fielddesc:LPTDATA1DATA}\item [data] 配置抽象 tdata1 的内容。由于仅支持匹配类型 (0x2) 触发器，此域将始终被解读为 \hyperref[regdesc:LPMCONTROL]{mcontrol} 的域。(R/W)
\end{reglist}\end{regdesc}
\end{register}


%%%%%%%%%
\begin{register}{H}{tdata2}{0x7A2}\label{regdesc:LPTDATA2}
\regfield{tdata2}{32}{0}{00000000000000000000000000000000}%
\reglabel{Reset}\regnewline%
\begin{regdesc}\begin{reglist}
\label{fielddesc:LPTDATA2}\item [tdata2] 配置抽象 tdata2 的内容。由于仅支持匹配类型 (0x2) 触发器，此域将始终被解读为 \hyperref[regdesc:LPMADDRESS]{maddress}。(R/W)
\end{reglist}\end{regdesc}
\end{register}


\begin{register}{H}{mcontrol}{0x7A1}\label{regdesc:LPMCONTROL}
\regfield{(reserved)}{4}{28}{{0x2}}%
\regfield{dmode}{1}{27}{{0}}%
\regfield{maskmax}{6}{21}{{0x0}}%
\regfield{hit}{1}{20}{{0}}%
\regfield{select}{1}{19}{{0}}%
\regfield{timing}{1}{18}{{0}}%
\regfield{sizelo}{2}{16}{{0}}%
\regfield{action}{4}{12}{{0001}}%
\regfield{chain}{1}{11}{{0}}%
\regfield{match}{4}{7}{{0}}%
\regfield{m}{1}{6}{{0}}%
\regfield{(reserved)}{1}{5}{{0}}%
\regfield{s}{1}{4}{{0}}%
\regfield{u}{1}{3}{{0}}%
\regfield{execute}{1}{2}{{0}}%
\regfield{store}{1}{1}{{0}}%
\regfield{load}{1}{0}{{0}}%
\reglabel{Reset}\regnewline%
\begin{regdesc}\begin{reglist}
\label{fielddesc:LPMCONTROLDMODE}\item [dmode] 与 \hyperref[regdesc:LPTDATA1]{tdata1} 的 \hyperref[fielddesc:LPTDATA1DMODE]{dmode} 一致。(RO)
\label{fielddesc:LPMCONTROLMASKMAX}\item [maskmax] 表示硬件支持的最大自然对齐的二次幂范围。\\
0：一个字节的范围，仅支持精确匹配\\
其他值：不支持\\
(RO)
\label{fielddesc:LPMCONTROLHIT}\item [hit] 硬件未实现，始终为 0。(RO)
\label{fielddesc:LPMCONTROLSELECT}\item [select] 选择地址匹配或数据匹配。\\
0：匹配虚拟地址\\
1：匹配读取或写入的数据值或执行的指令\\
注意：只实现了地址匹配，此域总为 0。\\
(RO)
\label{fielddesc:LPMCONTROLTIMING}\item [timing] 表示触发器何时被触发。\\
0：当匹配后，触发器会在指令执行前触发\\
1：当匹配后，触发器会在指令执行后触发\\
注意：此域总为 0。 \\
(RO)
\label{fielddesc:LPMCONTROLSIZELO}\item [sizelo] 仅支持任意大小匹配，此域总为 0。(RO)
\label{fielddesc:LPMCONTROLACTION}\item [action] 配置选定的触发器在触发时进行以下操作。\\
0x0：引起断点异常\\
0x1：进入调试模式（仅当 dmode = 1 时有效）\\
注意：仅支持进入调试模式，此域总为 1。\\
(RO)
\label{fielddesc:LPMCONTROLCHAIN}\item [CHAIN] 硬件未实现，始终为 0。(RO)
\label{fielddesc:LPMCONTROLMATCH}\item [match] 配置触发器进行数据/指令地址的匹配操作。\\
0x0：精确字节匹配，即与访问中某个字节对应的地址必须精确匹配\\ \hyperref[regdesc:LPMADDRESS]{maddress} 的值。
0x1：NAPOT 匹配，即访问中至少有一个字节处于  \hyperref[regdesc:LPMADDRESS]{maddress} 中规定的 NAPOT 区域。\\
注意：仅支持精确字节匹配，此域总为 0。\\
(R/W)
\label{fielddesc:LPMCONTROLM}\item [m] 置位使选定的触发器在机器模式下操作。(RO)
\label{fielddesc:LPMCONTROLS}\item [s] 置位使选定的触发器在监督模式下操作。不支持在监督模式下操作，此域总为 0。(RO)
\label{fielddesc:LPMCONTROLU}\item [u] 置位使选定的触发器在用户模式下操作。不支持在用户模式下操作，此域总为 0。(RO)
\item [见下页……]
\end{reglist}\end{regdesc}
\end{register}
\addtocounter{Regfloat}{-1}
\begin{register}{H}{mcontrol}{0x7A1}
\begin{regdesc}\begin{reglist}
\item [接上页……]
\label{fielddesc:LPMCONTROLEXECUTE}\item [execute] 置位使选定的触发器匹配指令的虚拟地址。(R/W)
\label{fielddesc:LPMCONTROLSTORE}\item [store] 置位使选定的触发器匹配存储器写操作的虚拟地址。硬件不支持，此域总为 0。(RO)
\label{fielddesc:LPMCONTROLLOAD}\item [load] 置位使选定的触发器匹配存储器读操作的虚拟地址。硬件不支持，此域总为 0。(RO)
\end{reglist}\end{regdesc}
\end{register}

\begin{register}{H}{maddress}{0x7A2}\label{regdesc:LPMADDRESS}
\regfield{maddress}{32}{0}{{0x00000000}}%
\reglabel{Reset}\regnewline%
\begin{regdesc}\begin{reglist}
\label{fielddesc:LPMADDRESS}\item [maddress] 配置选定的触发器执行匹配操作时使用的地址。(R/W)
\end{reglist}\end{regdesc}
\end{register}
